
\begin{boxK}{t-Test}

%\section{t-Test}
\begin{multicols}{2}

\subsection*{Mittelwert gegen Konstante} 
wobei c (constant) vorgegeben
$$t=\frac{\bar{x}-c}{\hat{\sigma}_{\bar{x}}}$$
$$\mathrm{df}=n-1$$

\subsection*{MWU, US} 

\begin{align*}
t&=\frac{\bar{x}_A-\bar{x}_B}{\hat{\sigma}_{\bar{x}_A-\bar{x}_B}}\\
\mathrm{df}&=(n_A-1)+(n_B-1)
\end{align*}
Hinweis: Berechnung der $\mathrm{df}$ gilt nur bei Gleichheit der Populationsvarianzen. 

\subsection*{MWU, AS}
$$t_{\mathrm{AS}}=\frac{\bar{D}}{\hat{\sigma}_{\bar{D}}}$$
$$\mathrm{df}_{\mathrm{AS}}=n-1$$
$D$: Differenz zwischen Mess-Werten

\columnbreak
\subsection*{$\alpha$-Fehler-Kumulierung}

Werden $j=\frac{k(k-1)}{2}$ paarweise Vergleiche durchgeführt, ergibt sich unter der Annahme unabhängiger Tests für die Wahrscheinlichkeit mindestens eines $\alpha$-Fehlers:

$$
\alpha_{\mathrm{kum}}
=1-(1-\alpha)^j
$$

Um die Familienfehlerwahrscheinlichkeit zu kontrollieren, ist eine Kontrastanalyse eine bessere Alternative zu paarweisen Vergleichen. Sind mehrere Tests zwingend erforderlich, werden häufig folgende Korrekturen verwendet:

\vspace{0.5em}
Bonferroni:
$$
\alpha'=\frac{\alpha}{j}
$$

Šidák:
$$
\alpha'=1-(1-\alpha)^{1/j}
$$

\end{multicols}
\end{boxK}
